\chapter{Reactor Control, Period, and Protection Systems}
\label{ch:control}

Stability theory tells us how a reactor responds to disturbances; control and
protection systems are how engineers ensure the reactor stays within the stable,
safe envelope and is shut down if it strays. This chapter connects the dynamics of
the previous chapters to the practical machinery of reactivity control, period
measurement, and emergency protection---and shows precisely which of these
safeguards were defeated or flawed at Chernobyl.

\section{Means of reactivity control}

Reactivity is adjusted by several means with very different speeds and purposes.
\textbf{Control rods}---neutron absorbers (boron carbide, hafnium, silver-indium-cadmium)
---provide fast, localised, and reliable control; they are the primary means of
prompt shutdown. \textbf{Soluble chemical shim}---boric acid dissolved in the
coolant of a PWR---provides slow, spatially uniform control of the large, slow
reactivity changes of fuel burnup and xenon. \textbf{Burnable poisons}---fixed
absorbers that deplete with irradiation---offset the excess reactivity of fresh
fuel. The art of control is to match each tool to the timescale of the reactivity
change it must counter.

\section{Rod worth and the integral/differential curves}

The reactivity inserted by a control rod depends nonlinearly on its position
because it is proportional to the local flux-squared (the rod removes neutrons in
proportion to the flux, and is most effective where the flux is highest---near the
core centre). The \emph{differential rod worth} $\dd\rho/\dd z$ is therefore
largest near mid-core and small at the ends, and the \emph{integral worth}
$\rho(z)$ is the characteristic S-shaped curve. Two consequences matter for
stability. First, a rod moving through the high-worth central region inserts
reactivity rapidly---a hazard if withdrawn carelessly. Second, the axial flux
shape determines where a rod's worth lies; in the RBMK, with the flux skewed to the
bottom of the core, the bottom-entering graphite displacers acted exactly where the
flux---and hence the harmful positive insertion---was concentrated.

\section{The reactor period and its measurement}

The reactor period $T=1/s$ (Chapter~\ref{ch:pk}) is the most direct measure of how
fast power is changing. Operators and protection systems monitor the period with
\emph{period meters} (or the equivalent ``startup rate'' in decades per minute). A
short period is an alarm: protection systems trip the reactor if the period falls
below a set limit (e.g.\ a few seconds), because a short period signals approach to
dangerous reactivity. The inhour equation links the period directly to reactivity,
so a period-meter is, in effect, a real-time reactivity gauge. Period-based trips
are a first line of defence against reactivity excursions---though against a
prompt-critical insertion, the period collapses faster than any trip can act, which
is why \emph{inherent} prompt feedback, not protection, is the true safeguard at
that extreme.

\section{Protection systems and defence in depth}

The reactor protection system (RPS) monitors many variables---neutron flux, rate
of change (period), coolant temperature, pressure, flow, and water level---and
initiates an automatic scram when any exceeds a safe limit. The design philosophy
is \emph{defence in depth}: multiple independent, redundant, and diverse barriers,
so that no single failure (and ideally no plausible combination) leads to harm.
A protection system is engineered to be \emph{fail-safe}: loss of power or signal
drives the rods in, not out.

\begin{remark}[Defeated protections at Chernobyl]
The Chernobyl test required, and the operators implemented, the deliberate
\emph{disabling} of several automatic protections---including the trip on turbine
stop and certain emergency-cooling and flux signals---so that the test would not be
interrupted. Defeating protections is sometimes necessary for testing, but it
removes the very defence-in-depth barriers that exist to catch the unexpected.
Combined with operation in an unstable regime, it left the reactor with neither
inherent stability nor engineered protection. The single remaining safeguard---the
AZ-5 scram---was itself flawed by the positive scram effect. Every layer of defence
had been removed or inverted.
\end{remark}

\section{The speed problem}

A recurring theme is the mismatch between the timescales of protection and of a
prompt excursion. Mechanical scram---unlatching rods and letting them fall or drive
in---takes on the order of one to a few seconds in a fast modern system, and was
about $18\unit{s}$ for full insertion in the RBMK. A delayed-critical transient
($\rho<\beta$) evolves over seconds and can be caught. A prompt-critical transient
($\rho>\beta$) evolves over milliseconds and cannot. This is why the design
objective is to make prompt criticality \emph{unreachable} by normal means and to
ensure that the prompt feedback is negative: protection systems guard the
delayed-critical regime, but only physics guards the prompt regime.

\keypoint{Control and protection systems keep a reactor within its stable envelope
and shut it down if it strays---but they act on the delayed-neutron timescale of
seconds. Against a prompt-critical excursion, only inherent prompt negative
feedback is fast enough. At Chernobyl the protections were disabled, the one
remaining scram was flawed, and the inherent feedback was positive: defence in
depth had been hollowed out at every layer.}

\section*{Exercises}
\addcontentsline{toc}{section}{Exercises}
\begin{enumerate}[leftmargin=*,itemsep=2pt]
\item Distinguish control rods, soluble shim, and burnable poisons by speed and purpose.
\item Sketch the integral and differential rod-worth curves and explain their shape.
\item Why does a short reactor period trigger a protective scram, and why is this useless against a prompt-critical insertion?
\item Define defence in depth and explain how it was hollowed out at Chernobyl.
\item Explain the fail-safe principle for a reactor protection system.
\end{enumerate}
